\documentclass[11pt]{article}
\usepackage{amsmath,amssymb,booktabs}
\usepackage[margin=1in]{geometry}
\title{Step 203: Diagonal Gram Correction}
\date{2026-05-16}
\begin{document}
\maketitle

\section*{Verdict}
\[
\boxed{\texttt{V\_diagonal\_corrected\_xi\_inconclusive}.}
\]

\section*{Burnol Kernel Convention}

The arXiv source for Burnol 2002 (math/0208121), equation 1, writes
\[
K(z_1,z_2)=
\frac{E(z_1)E(z_2)-E(1-z_1)E(1-z_2)}{z_1+z_2-1}.
\]
It then states that as a function of \(z_2\), this is the evaluator at
\(\overline{z_1}\), and as a function of \(z_1\), the evaluator at
\(\overline{z_2}\).  This is a bilinear de Branges normalization, not the
literal conjugated Hilbert-kernel substitution used naively in Step 202.

\section*{L'Hopital Diagonal}

Let \(w=1/2+i\tau\).  The norm of the evaluator at \(w\) is obtained by
taking \(z_1=\overline w\) and \(z_2\to w\):
\[
K(\overline w,w)
=\lim_{z\to w}
\frac{E(\overline w)E(z)-E(w)E(1-z)}{\overline w+z-1}.
\]
Differentiating the numerator with respect to \(z\) gives
\[
K(\overline w,w)
=E(\overline w)E'(w)+E(w)E'(\overline w).
\]
For the real entire \(E\)-function this is
\[
\boxed{K(\overline w,w)=2\operatorname{Re}\left(\overline{E(w)}E'(w)\right)}.
\]
The opposite sign was tested.  The plus sign is selected by the nonnegative
norm convention and is stable across the three tested zeros.

\section*{Derivative of \(E_{1/2}\)}

Burnol 2002, Theorem 8, gives
\[
E_\lambda(w)=A(w)B(w),\qquad
A(w)=\pi^{-w/2}\Gamma(w/2),
\]
\[
B(w)=\lambda^{1/2-w}
\frac{\sqrt\lambda}{2}\int_\lambda^\infty
(\psi_+^\lambda(t)-\psi_-^\lambda(t))t^{-w}\,dt.
\]
Therefore
\[
A'(w)=A(w)\left(-\frac{\log\pi}{2}+\frac{1}{2}\psi(w/2)\right),
\]
\[
B'(w)=
-\log(\lambda)\lambda^{1/2-w}
-\frac{\sqrt\lambda}{2}\int_\lambda^\infty
(\psi_+^\lambda(t)-\psi_-^\lambda(t))(\log t)t^{-w}\,dt.
\]
The script evaluates \(E'=A'B+AB'\) with the same finite cosine-resolvent
model used in Step 202.

\section*{Numerical Values}

\begin{center}
\begin{tabular}{c c c}
\toprule
zero & \(E'_{1/2}(\rho_i)\) & \(G_{ii}\)\\
\midrule
\(\rho_1\) & \(-2.18362493552\cdot10^{-5}+1.54944169115\cdot10^{-5}i\) & \(6.66157962829\cdot10^{-10}\)\\
\(\rho_2\) & \(-1.87039981590\cdot10^{-8}-1.01285320801\cdot10^{-7}i\) & \(1.28713995247\cdot10^{-14}\)\\
\(\rho_3\) & \(-4.34711344794\cdot10^{-9}-1.20194935955\cdot10^{-9}i\) & \(2.39433517220\cdot10^{-17}\)\\
\bottomrule
\end{tabular}
\end{center}

The displayed diagonal candidates are stable under \(N=160,240,320\) for the
finite resolvent and under \(50,70,100\) scalar decimal digits.  However, the
cutoff-derived error bounds are much larger than the values, so the numbers
are not certified.

\section*{Dual-System Cross-Check}

Burnol 2004 (math/0203120), Theorem 3.1, identifies the \(Y^a_{\rho,k}\)
minimal/complete regimes and describes the dual system.  Theorem 3.2 identifies
the \(Z^a_{\rho,k}\) minimal but not complete regime for \(a<1\).  The paper
also records the Euclid bilinear form \([f,g]=\int f(t)g(t)\,dt\).  These
records align structurally with the bilinear diagonal limit above, but this
step did not find an independent numerical norm formula for \(a=1/2\).

\section*{Full Gram and \(c\)-Matrix}

The full \(3\times3\) Gram matrix was assembled by using the L'Hopital diagonal
and Step 202 off-diagonal entries.  The commutator entries
\[
c_{ij}(\ell)=\langle e_i,(I-P_\infty)M_{m_\ell}P_\infty e_j\rangle
\]
were not certified, because the diagonal error dominates and the transported
\(\kappa\)-boundary samples plus the full \(K_\infty^{op}\) PSWF correction
still require a certified quadrature.

\section*{Nonclaim}

This step does not prove RH, does not close \(\Xi_{BC}\), and does not close
Branch B.  It fixes the symbolic diagonal and exposes the next numerical
certification gate.

\end{document}

